%-----------------------------------LICENSE------------------------------------%
%   This file is part of Mathematics-and-Physics.                              %
%                                                                              %
%   Mathematics-and-Physics is free software: you can redistribute it and/or   %
%   modify it under the terms of the GNU General Public License as             %
%   published by the Free Software Foundation, either version 3 of the         %
%   License, or (at your option) any later version.                            %
%                                                                              %
%   Mathematics-and-Physics is distributed in the hope that it will be useful, %
%   but WITHOUT ANY WARRANTY; without even the implied warranty of             %
%   MERCHANTABILITY or FITNESS FOR A PARTICULAR PURPOSE.  See the              %
%   GNU General Public License for more details.                               %
%                                                                              %
%   You should have received a copy of the GNU General Public License along    %
%   with Mathematics-and-Physics.  If not, see <https://www.gnu.org/licenses/>.%
%------------------------------------------------------------------------------%
\documentclass{article}
\usepackage{amssymb}
\usepackage{amsmath}

\title{18.100P: Midterm}
\date{Spring 2026}

% No indent and no paragraph skip.
\setlength{\parindent}{0em}
\setlength{\parskip}{0em}

\begin{document}
    \maketitle
    \textbf{Foundations / Logic}
    \begin{itemize}
        \item (6 Pts)
            Write down the truth table for:
            \begin{enumerate}
                \item
                    $P\Rightarrow\neg{Q}$
                \item
                    $Q\Rightarrow\neg{P}$
                \item
                    $(P\Rightarrow\neg{Q})\Rightarrow(Q\Rightarrow\neg{P})$
            \end{enumerate}
        \item (4 Pts)
            Is it consistent to accept this final expression
            as an axiom of logic? Explain using your truth table.
    \end{itemize}
    \newpage
    \textbf{Sequences and Subsequences}
    \begin{itemize}
        \item
            (2 Pts)
            Provide the precise definition of a convergent sequence of real
            numbers.
        \item
            (4 Pts)
            Prove that if $a,b:\mathbb{N}\rightarrow\mathbb{R}$ are
            convergent sequences with limits $A$ and $B$, respectively,
            then the sequence of complex numbers
            $z_{n}=a_{n}+ib_{n}$ also converges, and determine the limit.
            As a reminder, the distance between two complex numbers is given
            by the Pythagorean formula:
            \begin{equation}
                |(a+ib)-(c+id)|
                =
                \sqrt{(a-c)^{2}+(b-d)^{2}}
            \end{equation}
        \item
            (4 Pts)
            Prove or refute the claim that
            $a_{n}=\cos(n^{2})+\frac{1}{n+1}$ has a convergent subsequence
            (you may use any of the theorems from class without proving them).
    \end{itemize}
    \newpage
    \textbf{Cauchy Sequences}
    \begin{itemize}
        \item
            (2 Pts)
            Provide the precise definition of a Cauchy sequence of real numbers.
        \item
            (4 Pts)
            Prove that a convergent sequence of real numbers is also a
            Cauchy sequence.
        \item
            (4 Pts)
            Prove that if $a:\mathbb{N}\rightarrow\mathbb{R}$ is a Cauchy
            sequence, and if $f:\mathbb{R}\rightarrow\mathbb{R}$ is continuous
            everywhere, then $f(a_{n})$ also forms a Cauchy sequence.
    \end{itemize}
    \newpage
    \textbf{Compactness}
    \begin{itemize}
        \item
            (5 Pts)
            Prove that if $A\subseteq\mathbb{R}$ is \textit{complete}
            (meaning Cauchy sequences in $A$ also converge in $A$) and
            \textit{bounded}, then $A$ is compact
            (you may use either the subsequence definition or the open cover
            definition of compactness).
        \item
            (5 Pts)
            Prove that if $A\subseteq\mathbb{R}$ is compact, then it is
            complete (we proved in class that compact sets are bounded, so
            complete and bounded is an equivalent definition of compactness
            for $\mathbb{R}$).
            [Hint: You may use any of the major theorems from class freely.]
    \end{itemize}
    \newpage
    \textbf{Continuity}
    \begin{itemize}
        \item (10 Pts)
            Prove using an $\varepsilon-\delta$ argument that
            $f:\mathbb{R}\rightarrow\mathbb{R}$ given by:
            \begin{equation}
                f(x)
                =
                \begin{cases}
                    x^{2},&x\in\mathbb{Q}\\
                    -x^{2},&x\not\in\mathbb{Q}
                \end{cases}
            \end{equation}
            is continuous at $x=0$ and discontinuous everywhere else.
    \end{itemize}
\end{document}
